%%%%%%%%%%%%%%%%%%%%%%%%%%%%%%%%%%%%%%%%%%%%%%%%%%%%%%%%%%%%%%%%%%%%%
%% This is a (brief) model paper using the achemso class
%% The document class accepts keyval options, which should include
%% the target journal and optionally the manuscript type. 
%%%%%%%%%%%%%%%%%%%%%%%%%%%%%%%%%%%%%%%%%%%%%%%%%%%%%%%%%%%%%%%%%%%%%
\documentclass[journal=mamobx,manuscript=article]{achemso}

%%%%%%%%%%%%%%%%%%%%%%%%%%%%%%%%%%%%%%%%%%%%%%%%%%%%%%%%%%%%%%%%%%%%%
%% Place any additional packages needed here.  Only include packages
%% which are essential, to avoid problems later. Do NOT use any
%% packages which require e-TeX (for example etoolbox): the e-TeX
%% extensions are not currently available on the ACS conversion
%% servers.
%%%%%%%%%%%%%%%%%%%%%%%%%%%%%%%%%%%%%%%%%%%%%%%%%%%%%%%%%%%%%%%%%%%%%
\usepackage[version=3]{mhchem} % Formula subscripts using \ce{}

%%%%%%%%%%%%%%%%%%%%%%%%%%%%%%%%%%%%%%%%%%%%%%%%%%%%%%%%%%%%%%%%%%%%%
%% If issues arise when submitting your manuscript, you may want to
%% un-comment the next line.  This provides information on the
%% version of every file you have used.
%%%%%%%%%%%%%%%%%%%%%%%%%%%%%%%%%%%%%%%%%%%%%%%%%%%%%%%%%%%%%%%%%%%%%
%%\listfiles

%%%%%%%%%%%%%%%%%%%%%%%%%%%%%%%%%%%%%%%%%%%%%%%%%%%%%%%%%%%%%%%%%%%%%
%% Place any additional macros here.  Please use \newcommand* where
%% possible, and avoid layout-changing macros (which are not used
%% when typesetting).
%%%%%%%%%%%%%%%%%%%%%%%%%%%%%%%%%%%%%%%%%%%%%%%%%%%%%%%%%%%%%%%%%%%%%
% \usepackage{subfiles}
\usepackage{amssymb}
\usepackage{bm}
\usepackage{array}
\usepackage{booktabs}

% \usepackage{xr-hyper}
\usepackage[
  colorlinks=true,
  linkcolor=blue,
  citecolor=blue,
  urlcolor=blue,
  bookmarks=true,
  bookmarksopen=true,
  bookmarksnumbered=false,
  bookmarksdepth=2
]{hyperref}
\usepackage{xr}
\usepackage{silence}
\WarningFilter{hyperref}{Invalid page number}

\input{CrossPdfLinks.tex}
% The URL is evaluated from build/Manuscript.pdf, not from this source file.
\externaldocument[SI-]{./build/SupportingInformation}[SupportingInformation.pdf]

\usepackage{bookmark}
\bookmarksetup{
	open=true,
	depth=2
}

\newcounter{PdfSectionCounter}
\newcounter{PdfSubsectionCounter}

\newcommand{\PDFsection}[1]{%
  \stepcounter{PdfSectionCounter}%
  \hypertarget{pdfsection-\thePdfSectionCounter}{}%
  \bookmark[dest=pdfsection-\thePdfSectionCounter,level=1]{#1}%
  \section{#1}%
}

\newcommand{\PDFsubsection}[1]{%
  \stepcounter{PdfSubsectionCounter}%
  \hypertarget{pdfsubsection-\thePdfSubsectionCounter}{}%
  \bookmark[dest=pdfsubsection-\thePdfSubsectionCounter,level=2]{#1}%
  \subsection{#1}%
}

% 中文批注
% \usepackage[utf8]{inputenc}
% \usepackage{CJKutf8}
% \usepackage{xcolor}
% \usepackage[colorinlistoftodos,textsize=tiny]{todonotes}

% \newif\ifshowcomments
% % \showcommentstrue
% \showcommentsfalse

% \newcommand{\ctodo}[1]{%
%   \ifshowcomments
%     \todo[color=blue!15]{%
%       \begin{CJK*}{UTF8}{gbsn}%
%       \fontsize{7pt}{3pt}\selectfont #1%
%       \end{CJK*}%
%     }%
%   \fi
% }

\usepackage{lineno}

\newcommand*\mycommand[1]{\texttt{\emph{#1}}}

% ============================================================
% NT revision colors
% ============================================================

\usepackage{xcolor}
\usepackage{xparse}

% ---------- Color definitions ----------
\definecolor{ntBlue}{HTML}{326DB6}
\definecolor{ntRed}{HTML}{C44438}
\definecolor{ntOrange}{HTML}{F7842D}
\definecolor{ntGreen}{HTML}{4D7E54}
\definecolor{ntPurple}{HTML}{865A96}
\definecolor{ntTeal}{HTML}{008C8C}

% ---------- Line-number appearance ----------
\renewcommand{\linenumberfont}{%
  \normalfont
  \tiny
  \color{black}%
}

% ---------- Switch ----------
\newif\ifntshowcolor
% \ntshowcolortrue
\ntshowcolorfalse

% ---------- Core command ----------
\NewDocumentCommand{\ntc}{m m}{%
  \ifntshowcolor
    {\color{nt#1}#2}%
  \else
    #2%
  \fi
}

% ---------- Shortcuts ----------
\NewDocumentCommand{\ntr}{m}{\ntc{Red}{#1}}
\NewDocumentCommand{\ntb}{m}{\ntc{Blue}{#1}}
\NewDocumentCommand{\nto}{m}{\ntc{Orange}{#1}}
\NewDocumentCommand{\ntg}{m}{\ntc{Green}{#1}}
\NewDocumentCommand{\ntp}{m}{\ntc{Purple}{#1}}
\NewDocumentCommand{\ntt}{m}{\ntc{Teal}{#1}}

% ============================================================
% Revised text with automatic page and line labels
% ============================================================

% Label convention:
%   R1AQ1       Reviewer 1, group A, question 1
%   R1AQ1-a     First of several non-contiguous revisions
%   R2BQ3-b     Second revision for Reviewer 2, group B, question 3
%
% Usage:
% \ntq{R1AQ1-a}{
%   Revised manuscript text ...
% }
%
% Automatically creates:
%   rev:R1AQ1-a:start
%   rev:R1AQ1-a:end
%
% The labels remain active when revision colors are switched off.
\NewDocumentCommand{\ntq}{m +m}{%
  \leavevmode
  \phantomsection
  \label{rev:#1:StartAnchor}%
  \linelabel{rev:#1:start}%
  {\ntr{#2}}%
  \phantomsection
  \label{rev:#1:EndAnchor}%
  \linelabel{rev:#1:end}%
  \ignorespaces
}

% Location-only anchor: no color and no visible output.
% Use this at the beginning of a long or extensively reorganized passage.
% Usage: \ntqpoint{R1AQ1-b}
% Automatically creates: rev:R1AQ1-b:point
\NewDocumentCommand{\ntqpoint}{m}{%
  \phantomsection
  \label{rev:#1:PointAnchor}%
  \linelabel{rev:#1:point}%
}

%%%%%%%%%%%%%%%%%%%%%%%%%%%%%%%%%%%%%%%%%%%%%%%%%%%%%%%%%%%%%%%%%%%%%
%% Meta-data block
%% ---------------
%% Each author should be given as a separate \author command.
%%
%% Corresponding authors should have an e-mail given after the author
%% name as an \email command. Phone and fax numbers can be given
%% using \phone and \fax, respectively; this information is optional.
%%
%% The affiliation of authors is given after the authors; each
%% \affiliation command applies to all preceding authors not already
%% assigned an affiliation.
%%
%% The affiliation takes an option argument for the short name.  This
%% will typically be something like "University of Somewhere".
%%
%% The \altaffiliation macro should be used for new address, etc.
%% On the other hand, \alsoaffiliation is used on a per author basis
%% when authors are associated with multiple institutions.
%%%%%%%%%%%%%%%%%%%%%%%%%%%%%%%%%%%%%%%%%%%%%%%%%%%%%%%%%%%%%%%%%%%%%
\author{Shengda Zhao}
\author{Jiaxin Yu}
\affiliation[Beijing JiaoTong University]
{School of Physical Science and Engineering, Beijing JiaoTong University, Beijing 100044, PR China}
\author{Qiuju Chen}
\affiliation[Beijing Normal University]
{School of Physics and Astronomy, Beijing Normal University, Beijing 100875, PR China}
% \author{Yangjun Yan}
% \affiliation[Xihua University]
% {School of Science, Xihua University, Chengdu 610039, PR China}
\author{Quanxiao Dong}
\affiliation[Academy of Railway Sciences Corporation Limited Beijing]
{Railway Engineering Research Institute, China Academy of Railway Sciences Corporation Limited Beijing 100081, PR China}
\email{dongquanxiaotjs@163.com}
% \author{An-Chang Shi}
% \affiliation[McMaster University]
% {Department of Physics and Astronomy, McMaster University, Hamilton, Ontario L8S 4M1, Canada}
% \email{shi@mcmaster.ca}
\author{Xinghua Zhang}
\affiliation[Beihang University]
{School of Physics, Beihang University, Beijing, 100191, China}
\email{zhangxinghua@buaa.edu.cn}

%%%%%%%%%%%%%%%%%%%%%%%%%%%%%%%%%%%%%%%%%%%%%%%%%%%%%%%%%%%%%%%%%%%%%
%% The document title should be given as usual. Some journals require
%% a running title from the author: this should be supplied as an
%% optional argument to \title.
%%%%%%%%%%%%%%%%%%%%%%%%%%%%%%%%%%%%%%%%%%%%%%%%%%%%%%%%%%%%%%%%%%%%%
\title[Sliding Topology Regulates Microphase Separation]
  {Sliding Topology Regulates Microphase Separation in Linear-Star Topology-Transformable Copolymers}

%%%%%%%%%%%%%%%%%%%%%%%%%%%%%%%%%%%%%%%%%%%%%%%%%%%%%%%%%%%%%%%%%%%%%
%% Some journals require a list of abbreviations or keywords to be
%% supplied. These should be set up here, and will be printed after
%% the title and author information, if needed.
%%%%%%%%%%%%%%%%%%%%%%%%%%%%%%%%%%%%%%%%%%%%%%%%%%%%%%%%%%%%%%%%%%%%%
\abbreviations{SCFT, ODT, BCC, FCC}
\keywords{topology-transformable copolymers, sliding topology, self-consistent field theory, microphase separation, packing frustration}

%%%%%%%%%%%%%%%%%%%%%%%%%%%%%%%%%%%%%%%%%%%%%%%%%%%%%%%%%%%%%%%%%%%%%
%% The manuscript does not need to include \maketitle, which is
%% executed automatically.
%%%%%%%%%%%%%%%%%%%%%%%%%%%%%%%%%%%%%%%%%%%%%%%%%%%%%%%%%%%%%%%%%%%%%
\begin{document}

%%%%%%%%%%%%%%%%%%%%%%%%%%%%%%%%%%%%%%%%%%%%%%%%%%%%%%%%%%%%%%%%%%%%%
%% The "tocentry" environment can be used to create an entry for the
%% graphical table of contents. It is given here as some journals
%% require that it is printed as part of the abstract page. It will
%% be automatically moved as appropriate.
%%%%%%%%%%%%%%%%%%%%%%%%%%%%%%%%%%%%%%%%%%%%%%%%%%%%%%%%%%%%%%%%%%%%%
\begin{tocentry}
\centering
\includegraphics[]{./Figure/TOC.png}
% Some journals require a graphical entry for the Table of Contents.
% This should be laid out ``print ready'' so that the sizing of the
% text is correct.

% Inside the \texttt{tocentry} environment, the font used is Helvetica
% 8\,pt, as required by \emph{Journal of the American Chemical
% Society}.

% The surrounding frame is 9\,cm by 3.5\,cm, which is the maximum
% permitted for  \emph{Journal of the American Chemical Society}
% graphical table of content entries. The box will not resize if the
% content is too big: instead it will overflow the edge of the box.

% This box and the associated title will always be printed on a
% separate page at the end of the document.

% 8.25 cm × 4.45 cm。
\end{tocentry}

%%%%%%%%%%%%%%%%%%%%%%%%%%%%%%%%%%%%%%%%%%%%%%%%%%%%%%%%%%%%%%%%%%%%%
%% The abstract environment will automatically gobble the contents
%% if an abstract is not used by the target journal.
%%%%%%%%%%%%%%%%%%%%%%%%%%%%%%%%%%%%%%%%%%%%%%%%%%%%%%%%%%%%%%%%%%%%%
\begin{abstract}
	Self-assembled polymeric structures are often designed by predetermined molecular parameters, and studies that directly connect reconfigurable molecular topology to microphase-separated morphologies remain limited.
	Here, we develop a self-consistent field theory (SCFT) for linear--star topology-transformable (L--S TT) copolymers, in which the grafting position is treated as a continuous sliding degree of freedom, $f$.
   Numerical solutions of the SCFT equations yield a globally preferred topology, $f_\Lambda^\ast$, consistent with the most unstable fixed-topology RPA mode near the spinodal.
  The corresponding RPA spinodal difference defines an effective conformational-entropy scale, $\Delta S_{\mathrm{conf}}$, while normalized local topology distributions quantify the sliding-distribution entropy loss, $\Delta S_{\mathrm{Slide}}$, relative to uniform sliding.
  Moreover, the spatial redistribution of sliding configurations is obtained, revealing local topological adaptation to interfacial geometry and packing constraints.
  The predicted phase diagrams further demonstrate that sliding topology enhances the stability of spherical phases.
  More broadly, the theory can be generalized to diverse sliding-topology copolymers and provides a transferable quantitative framework for probing topology redistribution in response to curvature and packing constraints, offering guidance for the design of complex ordered phases and micelles in copolymer systems.
\end{abstract}

% TODO:关键词
%%%%%%%%%%%%%%%%%%%%%%%%%%%%%%%%%%%%%%%%%%%%%%%%%%%%%%%%%%%%%%%%%%%%%
%% Start the main part of the manuscript here.
%%%%%%%%%%%%%%%%%%%%%%%%%%%%%%%%%%%%%%%%%%%%%%%%%%%%%%%%%%%%%%%%%%%%%

\linenumbers

\PDFsection{Introduction}
	\input{sections/Introduction.tex}

\PDFsection{Theory and Methods}
	\input{sections/Methods.tex}

\PDFsection{Numerical Methods}
	\input{sections/NumericalMethods.tex}

\PDFsection{Results and Discussion}
	\input{sections/ResultsOverview.tex}

\PDFsection{Conclusion}
    \input{sections/Conclusion.tex}
%%%%%%%%%%%%%%%%%%%%%%%%%%%%%%%%%%%%%%%%%%%%%%%%%%%%%%%%%%%%%%%%%%%%%
%% The "Acknowledgement" section can be given in all manuscript
%% classes.  This should be given within the "acknowledgement"
%% environment, which will make the correct section or running title.
%%%%%%%%%%%%%%%%%%%%%%%%%%%%%%%%%%%%%%%%%%%%%%%%%%%%%%%%%%%%%%%%%%%%%
\begin{acknowledgement}
	This work was supported by the National Natural Science Foundation of China (NSFC) 
	(Nos. 22573005 and U22A20150) 
	and the Fundamental Research Funds for the Central Universities 
	(No. 2025JBZY027).
    % This work was supported by 
    % National Natural Science Foundation of China (NSFC) No. 22573005, 
    % the Fundamental Research Funds for the Central Universities No. 2025JBZY027
	% and National Natural Science Foundation of China (NSFC) No. U22A20150.

\end{acknowledgement}

%%%%%%%%%%%%%%%%%%%%%%%%%%%%%%%%%%%%%%%%%%%%%%%%%%%%%%%%%%%%%%%%%%%%%
%% The same is true for Supporting Information, which should use the
%% suppinfo environment.
%%%%%%%%%%%%%%%%%%%%%%%%%%%%%%%%%%%%%%%%%%%%%%%%%%%%%%%%%%%%%%%%%%%%%
% \begin{suppinfo}
% \end{suppinfo}

%%%%%%%%%%%%%%%%%%%%%%%%%%%%%%%%%%%%%%%%%%%%%%%%%%%%%%%%%%%%%%%%%%%%%
%% The appropriate \bibliography command should be placed here.
%% Notice that the class file automatically sets \bibliographystyle
%% and also names the section correctly.
%%%%%%%%%%%%%%%%%%%%%%%%%%%%%%%%%%%%%%%%%%%%%%%%%%%%%%%%%%%%%%%%%%%%%
% \bibliography{achemso-demo}
% \bibliography{C:/Users/79323/OneDrive/Working/DailyRead/biblib}
\bibliography{cleaned}

\end{document}
